\chapter{分离变量法}

\problem{14.1}{长为 $l$ 的弦，两端固定，弦中的张力为 $T$，在距弦的一端 $x=0$ 的 $l_0$ 处（$l_0<l$），以力 $F_0$ 把弦拉开，使 $l_0$ 点离开平衡位置有一小位移，并使弦处于稳定状态。在 $t=0$ 时刻突然撤销这一大小为 $F_0$ 的力，求解弦的振动。}
\answerarea

\problem{14.2}{长为 $l$ 的弦，两端固定，在弦的 $x_0$ 处（$0<x_0<l$）施予谐变外力 $\rho f(t)=\rho A\sin\omega t$，（$\rho$ 为弦的线密度，$A$ 为常数），求解弦的振动。

（提示：外力可表示为 $\rho f(x,t)=\rho A\delta(x-x_0)\sin\omega t$。）}
\answerarea
\clearpage

\problem{14.3}{求解习题 13.1 中的（2）、（3）两种情况下杆的振动。}
\answerarea

\problem{14.4}{导热系数为 $\kappa$ 的均匀物质构成的立方体，$0\leqslant x\leqslant l_1,\ 0\leqslant y\leqslant l_2,\ 0\leqslant z\leqslant l_3$，其初始温度为 $T_0$，如果其表面温度保持在 $0\,^\circ\mathrm{C}$，求此立方体温度的变化。并写出此立方体中心点，即 $\left(\dfrac{l_1}{2},\dfrac{l_2}{2},\dfrac{l_3}{2}\right)$ 处温度随时间的变化关系。又问当时间 $t\to\infty$ 时，该立方体的温度表达式，并从物理上解释其合理性。

思考：如果边界条件改为该立方体置于 $0\,^\circ\mathrm{C}$ 的环境中，其边界以牛顿冷却定律与环境进行热交换，其解的形式又是如何变化的，并与上面的结果进行比较，分析这两种边界条件的物理图像上的差异。}
\answerarea
\clearpage

\problem{14.5}{侧面绝热的均匀细杆，$0\leqslant x\leqslant l$，其导热系数为 $\kappa$，初始温度为 $T_0$，其两端保持不变的温度：$u(0,t)=T_1,\,u(l,t)=T_2$，$T_1,T_2$ 均为常数。求杆上的温度分布随时间的变化关系。当 $t\to\infty$ 时，给出杆上的温度分布的表达式，并解释这一结果在物理上的合理性及其满足的物理定律。}
\answerarea

\problem{14.6}{铀块中中子的扩散、吸收的增殖反应如习题 13.4 所述，试求一维细铀杆的临界长度。临界长度是指当铀杆超过这一长度时，杆中的中子浓度随时间而增长致使铀块产生爆炸。这里忽略杆的过中心轴剖面的对角线长度与杆的长度之差。}
\answerarea
\clearpage

\problem{14.7}{一厚度为 $l$ 的单晶硅片置于充满杂质的环境中，杂质从硅片的两表面向里扩散，假设硅片两表面杂质浓度保持恒定的 $N_0$，杂质在硅片中的扩散系数为 $D$，求解此扩散问题。并说明扩散的时间与扩散系数 $D$ 及厚度 $l$ 的关系。当 $t$ 较大时，化简所得的结果。}
\answerarea

\problem{14.8}{在矩形区域 $0<x<a,\ 0<y<b$ 上求解 Laplace 方程
\[
  \nabla^2 u(x,y) = 0
\]
使其满足边界条件：
\[
  \left\{\begin{aligned}
    u(0,y) &= Ay(b-y), & u(a,y) &= 0 \\
    u(x,0) &= B\sin\frac{\pi x}{a}, & u(x,b) &= 0
  \end{aligned}\right.
\]
其中 $A,B$ 为常数。}
\answerarea
\clearpage

\problem{14.9}{在量子物理学中，空间中自由粒子可视为 de Broglie（德布罗意）波，被称为物质波，用 $\psi(\boldsymbol{r},t)$ 来描述，它满足如下 Schrödinger（薛定谔）方程（量子波动方程）
\[
  \mathrm{i}\hbar\frac{\partial}{\partial t}\psi(\boldsymbol{r},t) = -\frac{\hbar^2}{2m}\nabla^2\psi(\boldsymbol{r},t)
\]
其中 $\hbar$ 为 Planck（普朗克）常数，$m$ 为粒子质量。我们把粒子看作在一个三边长为 $L$ 的正立方体 $\left(-\dfrac{L}{2}\leqslant x\leqslant\dfrac{L}{2},\ -\dfrac{L}{2}\leqslant y\leqslant\dfrac{L}{2},\ -\dfrac{L}{2}\leqslant z\leqslant\dfrac{L}{2}\right)$ 中的三维波动，且这一波动满足周期性边界条件，试求解这一波动。}
\answerarea
\clearpage
